% ============================
% パッケージ
% ============================
\usepackage{amsmath,amssymb,amsthm}
\usepackage{mathrsfs}
\usepackage{tabularx,longtable,multirow,here,colortbl}
\definecolor{pink}{cmyk}{0.093, 0.551, 0.289, 0}
\usepackage[dvipdfmx]{graphicx}
\usepackage{wrapfig}
\usepackage{ascmac,comment,url,makeidx,subfiles,pifont}
\usepackage{tcolorbox}
\tcbuselibrary{breakable, skins, theorems}
\usepackage{enumitem}
\usepackage{tikz}
\usetikzlibrary{intersections, calc, arrows, trees}
\usepackage[dvipdfmx,bookmarkstype=toc,colorlinks=true,urlcolor=black,linkcolor=blue,citecolor=blue,linktocpage=true,bookmarks=true]{hyperref}
\AtBeginDvi{\special{pdf:tounicode 90ms-RKSJ-UCS2}}
\usepackage{geometry}
\geometry{top=2cm, bottom=2cm, left=1cm, right=1cm, includefoot}

% ============================
% 定理環境
% ============================
\newtheoremstyle{mystyle}{5pt}{12pt}{\normalfont\setlength{\parindent}{2em}}{}{\bf}{.}{\newline}{\underline{\thmname{#1}{\;}\thmnumber{#2}}\thmnote{（#3）}}
\theoremstyle{mystyle}
\newtheorem{my-theorem}{Theorem}[section]
\newtheorem{my-lemma}[my-theorem]{Lemma}
\newtheorem{my-proposition}[my-theorem]{Proposition}
\newtheorem{my-corollary}[my-theorem]{Corollary}
\newtheorem{my-definition}[my-theorem]{Definition}
\newtheorem{my-example}[my-theorem]{Example}
\newtheorem{my-remark}[my-theorem]{Remark}
\newtheorem{my-exercise}[my-theorem]{Exercise}

\newcommand{\thmsymbol}{$\blacksquare$}

% 環境上書き
\newenvironment{theorem}[1][] {\begin{my-theorem}[#1]\renewcommand{\qedsymbol}{\thmsymbol}\pushQED{\qed}}{\popQED\end{my-theorem}}
\newenvironment{lemma}[1][] {\begin{my-lemma}[#1]\renewcommand{\qedsymbol}{\thmsymbol}\pushQED{\qed}}{\popQED\end{my-lemma}}
\newenvironment{proposition}[1][] {\begin{my-proposition}[#1]\renewcommand{\qedsymbol}{\thmsymbol}\pushQED{\qed}}{\popQED\end{my-proposition}}
\newenvironment{corollary}[1][] {\begin{my-corollary}[#1]\renewcommand{\qedsymbol}{\thmsymbol}\pushQED{\qed}}{\popQED\end{my-corollary}}
\newenvironment{definition}[1][] {\begin{my-definition}[#1]\renewcommand{\qedsymbol}{\thmsymbol}\pushQED{\qed}}{\popQED\end{my-definition}}
\newenvironment{example}[1][] {\begin{my-example}[#1]\renewcommand{\qedsymbol}{\thmsymbol}\pushQED{\qed}}{\popQED\end{my-example}}
\newenvironment{remark}[1][] {\begin{my-remark}[#1]\renewcommand{\qedsymbol}{\thmsymbol}\pushQED{\qed}}{\popQED\end{my-remark}}
\newenvironment{exercise}[1][] {\begin{my-exercise}[#1]\renewcommand{\qedsymbol}{\thmsymbol}\pushQED{\qed}}{\popQED\end{my-exercise}}

% ============================
% 証明環境
% ============================
\renewcommand{\proofname}{Proof}
\makeatletter
\renewenvironment{proof}[1][\proofname]{\par\pushQED{\qed}\normalfont \topsep6\p@\@plus6\p@\relax \trivlist\item[\hskip\labelsep\itshape {\bf\underline{#1}\@addpunct{\ }}]\ignorespaces}{\popQED\endtrivlist\@endpefalse}
\makeatother
\newenvironment{answer}{\begin{proof}[Answer]}{\end{proof}}

% ============================
% マクロ
% ============================
\newcommand{\defarr}{\ \mbox{$\stackrel{\text{def}}{\Longleftrightarrow}$}}
\newcommand{\defeq}{\ \mbox{$\stackrel{\text{def}}{=}$}}
\newcommand{\tempdefeq}{:=}
\newcommand{\map}[3]{#1\colon #2 \to #3}
\newcommand{\InjectionMap}[3]{#1\colon #2 \xrightarrow{\text{inj}} #3}
\newcommand{\SurjectionMap}[3]{#1\colon #2 \xrightarrow{\text{sur}} #3}
\newcommand{\BijectionMap}[3]{#1\colon #2 \xrightarrow{\text{bij}} #3}
\newcommand{\op}[2]{\langle #1,#2 \rangle}
\DeclareMathOperator{\dom}{dom}
\DeclareMathOperator{\ran}{ran}
\newcommand{\idmap}[1]{\mathrm{id}_{#1}}
\DeclareMathOperator{\card}{card}
\newcommand{\SetBar}[2]{\{\,#1\,\mid\,#2\,\}}
\newcommand{\tempdefarr}{:\Leftrightarrow}
\newcommand{\dash}[1]{#1 ^{\prime}}
\newcommand{\ImpDir}[2]{\hspace{-0.5cm}（\;(#1)$\Rightarrow$(#2)\;）}

% sub-block 環境
\definecolor{BIP-back}{gray}{0.85}
\newcount\subblocklevel
\subblocklevel=0
\newcommand{\subblockcolback}{\ifodd\subblocklevel BIP-back\else white\fi}
\makeatletter
\newcommand{\sbfootnote}[1]{\footnotemark\expandafter\g@addto@macro\csname sb@footnotes@\the\subblocklevel\endcsname{\footnotetext{#1}}}
\makeatother
\newenvironment{sub-block}[1]{\global\advance\subblocklevel by 1\relax\expandafter\gdef\csname sb@footnotes@\the\subblocklevel\endcsname{}\vspace{-5pt}\begin{tcolorbox}[colback = \subblockcolback,colframe = white,boxrule = 5pt,fonttitle = \bfseries,breakable = true]\begin{itemize} \item[#1]}{\end{itemize}\end{tcolorbox}\vspace{-5pt}\csname sb@footnotes@\the\subblocklevel\endcsname\global\advance\subblocklevel by -1\relax}
